\section{Step 14 Add choice, old h, the constant and the schedule word}\label{sha-s14}

The factored choice expression selects between old f and old g using old e. Add it to the capital-sigma-one result, then include old h, the current constant and the current schedule word. Use nontrapping word additions throughout. \texttt{t0} is now the complete T1; the state and working registers are still unchanged.

\subsection{Write and inspect}
Replace only the temporary instructions between \texttt{\# BEGIN comp\_t1} and \texttt{\# END comp\_t1}. Keep the supplied labels outside that region.

\Needspace{13\baselineskip}
\MARSCodeLabel{MIPS assembly}
\begin{minted}[samepage=true]{mars-idea}
    xor $t2, $s5, $s6
    and $t2, $t2, $s4
    xor $t2, $t2, $s6
    addu $t0, $t0, $t2
    addu $t0, $t0, $s7
    lw $t2, 0($t7)
    addu $t0, $t0, $t2
    lw $t2, 0($t6)
    addu $t0, $t0, $t2
\end{minted}

Use recipe \texttt{s14-1} with the cumulative program through this edit. The exact recipe names the entry, prepared input and checkpoint. Observe the checkpoint before continuing to its temporary or normal return.

\subsection{Expected observation}
For the first abc round, \texttt{t0} is \texttt{54da50e8}.

\Needspace{9\baselineskip}
\subsection{Check your understanding}
What does a one bit of e select in choice?

\begin{itemize}
  \item \textbf{A} The corresponding bit of f.
  \item \textbf{B} The corresponding bit of g.
  \item \textbf{C} The next constant.
\end{itemize}
